\begin{center}\begin{tikzpicture}[font=\small,>=stealth,line join=round]
\draw[gray!40,rounded corners=3pt] (0,0) rectangle (14.2,8.5);\node at (7.1,7.9) {例三：先确定原点，再用给定函数值求参数};
\begin{scope}[shift={(4.2,3.8)},x=1.7cm,y=0.85cm]\draw[->,gray] (-1.8,0)--(1.8,0) node[right]{$x$};\draw[->,gray] (0,-3.1)--(0,3.1) node[above]{$y$};\draw[blue,thick,domain=-1.1:1.1,samples=150] plot(\x,{(\x)^3+\x});\fill[orange!85!black] (1,2) circle(2.2pt) node[above right=5pt] {$(1,2)$};\fill[orange!85!black] (-1,-2) circle(2.2pt) node[below left=5pt] {$(-1,-2)$};\fill[orange!85!black] (0,0) circle(2.2pt) node[below right=5pt] {$O$};\end{scope}\node[align=left,anchor=west] at (8.3,4.1){求参后 $f(x)=x^3+x$\\[11pt]$a=c=0,\ b=1$\\[11pt]已知 $f(1)=2$\\[11pt]同时必有 $f(-1)=-2$};
\node[text width=12.8cm,align=center] at (7.1,.65) {图像画的是求参后的实际函数；正负两点关于原点对称。};
\end{tikzpicture}\end{center}
